
\section{The Arena}\label{sec:arch-arena}

\subsection{What the Arena Is}\label{sec:arch-overview}

\textbf{The arena does not query models.} It re-judges responses that
were generated and stored earlier, in an evaluation corpus that also
records each model's selected answer. A run therefore consists of
routing held-out queries through a frozen taxonomy snapshot, retrieving
two stored responses per comparison, and issuing judge calls. No
response is produced during a run.

This is stated first because it is the single largest difference
between this system and the way an ``arena'' is usually described, and
because two later results are unintelligible without it. First, a model
whose response text was never persisted cannot be distinguished
downstream from one whose text was: the trace accessor never returns an
empty string, and synthesises a one-line stub from the stored answer
instead. That is the mechanism behind the capture-gap defect of
Section~\ref{sec:res-capture-gap}. Second, because the responses are
fixed, the two judging conditions of Chapter~\ref{ch:experimental-design}
see byte-identical inputs, which is what makes the comparison paired.

\paragraph{Notation.}
Let $\mathcal{Q}$ be the corpus of MMLU-Pro queries, partitioned once
before construction into a construction set
$\mathcal{Q}_{\text{con}}$ and a disjoint evaluation set
$\mathcal{Q}_{\text{eval}}$.
$\mathcal{Q}_{\text{con}}$ is the only query set visible to taxonomy
construction and to rubric generation; $\mathcal{Q}_{\text{eval}}$ is
reserved for held-out routing and judging.
At depth $\ell$ the effective representation of $q_i$ is a fixed
256-dimensional prefix slice
$\mathbf{x}_i^{(\ell)} \in \mathcal{S}^{d-1}$.
The anchor set $\mathcal{A} = \{D_1, \ldots, D_{14}\}$ is in one-to-one
correspondence with the MMLU-Pro domains; each anchor carries a tree of
internal nodes and leaves $\mathcal{L}(D_k)$, and because every node
has exactly one parent the sets $\mathcal{L}(D_k)$ partition the leaf
set $\mathcal{L}$.
The \textit{canonical partition}
$Q_{D_k}^{\text{canonical}} = \{q_i : y_i^{\text{MMLU}} = D_k\}$ is the
editorial one; the \textit{adapted partition}
$Q_{D_k}^{\text{adapted}} = \{q_i : \hat{D}(q_i) = D_k\}$ is the one
construction induces, where $\hat{D}(q_i)$ is the anchor of $q_i$'s
primary leaf. The two coincide absent cross-domain migration.
For a set of evaluated models $\mathcal{M} = \{m_1, \ldots, m_M\}$ the
estimand is a per-cell Bradley--Terry log-strength vector
$\boldsymbol{\theta}$, and the validation objective is its rank
agreement with per-cell ground-truth accuracy.

\subsection{Judge Design}\label{sec:arch-judge}

Each leaf carries a rubric induced once per frozen snapshot. The
rubric is generated by a MapReduce procedure over the leaf's
construction-pool queries: packs of $k = 25$ queries are presented to a
strong LLM together with their MMLU-Pro reference answers and prompted
for the quality principles that separate correct from incorrect
responses for queries of that type; a reduce call consolidates the
per-pack principle lists into one ranked, leaf-specific rubric.
Rubrics are stored with the snapshot, so every arena run against a
snapshot uses identical criteria.

\textbf{Held-out queries never enter rubric generation, and this is
audited rather than asserted.} The induction pool is filtered against
the reserved set, and the filter is accompanied by an assertion that
recomputes the membership test independently and \emph{fails the build}
if the two disagree. This is the strongest information-separation
guarantee in the system. A further detail belongs with it: the
induction pool is additionally thinned by roughly a fifth by a hash
filter, so a leaf's rubric is induced from a sample of its construction
queries rather than all of them.

At judging time each comparison $(m_i, m_j, q, L)$ is presented to the
judge as a system prompt carrying leaf $L$'s rubric, followed by a
question block and the two responses. The shared structural template
covering presentation, tie handling and output format is
\emph{appended} after the leaf-specific criteria, not prepended.
Each comparison is issued twice with the presentation order swapped; a
verdict counts as consistent when both calls name the same winner, and
a disagreement is recorded as \texttt{TIE}. This converts positional
disagreement into an explicit tie rather than leaving it as noise in
the sufficient statistics.

\textbf{What the judge does and does not see.}
The judge is \emph{answer-key-blind}: the prompt-construction path
asserts that the rendered query does not contain a ground-truth answer
block, and the stored correct answer is never read on the judging path.
The judge is \emph{not} answer-blind. The question block includes the
full multiple-choice options, and $94.8\%$ of the stored traces state
the answering model's own selected option. A judge capable of solving
the item can therefore verify both responses against its own solution
without ever being told the key. Chapter~\ref{ch:results} reports that
this is what it does, and Chapter~\ref{ch:discussion} treats it as a
boundary condition on the architecture rather than as a defect of it.

Four further properties of the judging path affect how verdicts should
be read.

\begin{enumerate}
\item The verdict schema carries a numeric \texttt{confidence}
alongside the winner, and that field is consumed downstream rather than
being decorative.
\item Consolidated confidence is forced to $0.5$ whenever the two
presentation orders disagree, and is capped at $0.95$ otherwise.
\item A post-hoc adjustment rewrites a pair's win counts when the
order asymmetry of its verdicts exceeds $0.3$ with at least six
comparisons. This is a correction applied to the sufficient statistics,
not to the verdicts, and it is applied before any fit.
\item A verdict that cannot be parsed is dropped silently: there is no
counter and no export row for it, so the realised comparison count is a
lower bound on the calls issued.
\end{enumerate}

\subsection{Bradley--Terry Estimation}\label{sec:meth-bt}

The Bradley--Terry model assigns each model $m_i$ a latent strength
$s_i$ with $P(m_i \succ m_j) = \sigma(s_i - s_j)$
\parencite{bradley1952rank}, and a zero-sum constraint
$\sum_i s_i = 0$ fixes the additive degree of freedom. Ties contribute
a half-win to each side, so $W_{ij} + W_{ji} = N_{ij}$ without an
explicit draw parameter; no Davidson tie-propensity term is fitted
\parencite{davidson1970ties}.

The estimator is the Minorisation--Maximisation update of
\textcite{hunter2004mm}, written in the exponentiated parameters
$\pi_i = e^{s_i}$:
\[
\pi_i^{(t+1)} \;=\;
\frac{W_i}{\displaystyle\sum_{j \neq i}
\frac{N_{ij}}{\pi_i^{(t)} + \pi_j^{(t)}}},
\qquad W_i = \sum_{j \neq i} W_{ij},
\]
applied as a simultaneous (Jacobi) sweep over all models, with the
zero-sum constraint re-enforced after each sweep. Iteration stops at a
tolerance of $\max(10^{-6},\, 0.01\,\widehat{\mathrm{SE}}_{\text{lb}})$
or at a cap of 200 sweeps.

A Jeffreys $\mathrm{Beta}(0.5, 0.5)$ prior is applied as a phantom
half-win in each direction, and only to pairs that already carry at
least one real comparison. Pairs with no data are left untouched, so
the prior regularises estimates rather than manufacturing edges in the
comparison graph.

\textbf{Standard errors.}
$\hat{\sigma}_i = \sqrt{[\mathcal{I}(\hat{\mathbf{s}})^{-1}]_{ii}}$ is
taken from the full observed Fisher information. The zero-sum
constraint makes $\mathcal{I}$ rank-deficient by one, and the
implementation does not use a pseudoinverse: it inverts the
rank-completed matrix $\mathcal{I} + \mathbf{1}\mathbf{1}^{\!\top}$ and
then removes the completion term. A hard floor of $0.01$ is applied to
every standard error.

\textbf{A correction to every standard error reported by earlier
versions of this pipeline.}
The rank completion requires subtracting $1/K^{2}$ after inversion,
not $1/K$; the implementation subtracted $1/K$, so \emph{every
Bradley--Terry standard error this system reported before the
correction was the floor constant rather than a measurement}. No
confidence interval computed before the correction is quoted in this
thesis. The derivation is in Appendix~\ref{app:numerics}; the
correction and its consequences are recorded once, in the corrections
register (Appendix~\ref{app:corrections-register}).

\textbf{Identifiability.}
Bradley--Terry parameters are identified only up to an additive
constant per connected component of the comparison graph, so a leaf
whose graph is disconnected yields incomparable sub-scales. The
scheduler opens with a connectivity bootstrap, but that bootstrap
operates on \emph{global} pair counts rather than per-leaf ones. It
therefore does not guarantee that every leaf's comparison graph is
connected, and the word ``guarantees'' is not used of it here.

\subsection{Scheduling and Stopping}\label{sec:meth-scheduling}

The scheduler assigns each candidate comparison triple $(i, j, L)$ a
composite utility
\[
U(i,j,L) = \left[\alpha\, U_{\text{struct}}(i,j,L)
+ (1-\alpha)\, U_{\text{resolve}}(i,j,L)\right]
\times \left(1 - 0.3\, \frac{n_{ij}^{(L)}}{B_{ij}}\right),
\]
and selects the highest-scoring batch.
$U_{\text{struct}}$ rewards closely ranked pairs; $U_{\text{resolve}}$
is the binary entropy of the current win probability
$p_{ij} = \sigma(\hat{s}_i - \hat{s}_j)$ divided by the combined
variance $\hat{\sigma}_i^2 + \hat{\sigma}_j^2$ and weighted by a KDE
valley weight favouring midpoints in low-density regions of the score
distribution (Appendix~\ref{app:numerics}).
The entropy term is maximised at $p_{ij} \approx 0.5$; the variance
division downweights pairs whose scores are already noisy, and
correspondingly upweights pairs that are ambiguous \emph{despite}
tight estimates --- the pairs where one more comparison is most likely
to be decisive.
The mixing parameter $\alpha$ is maturity-driven and falls as a leaf's
average standard error falls, so scheduling is exploration-dominated
early and resolution-dominated late. Two further terms are applied and
are recorded here for completeness: a mask that suppresses pairs whose
recent verdicts have been ties, and the repeat discount already visible
in the trailing factor above.

A pair resolves when
$|\hat{s}_i - \hat{s}_j| \geq 3.0(\hat{\sigma}_i + \hat{\sigma}_j)$ or
when its per-pair budget is exhausted. A run ends when the round
counter reaches its maximum; that is the de-facto exit in practice.

Two properties of the scheduler bear directly on how results must be
read. First, \textbf{it over-samples uncertain pairs}, so mid-ranked
models face harder average opposition than a uniform design would give
them. Win-rate ignores opponent strength and is therefore biased by
this; Bradley--Terry $\theta$ is not, and Chapter~\ref{ch:results}
quotes $\theta$ throughout; the measured cost of the difference is
reported with the result it affects (Section~\ref{sec:res-c5}).
Second, \textbf{query selection within a scheduled pair is a persisted
round-robin offset over the leaf's pool}, not a descending-cosine walk
from the leaf centroid. An earlier centrality-ordered sampler consumed
only about two-thirds of each leaf's pool and systematically omitted the
least prototypical third, so runs under the two samplers are not
comparable; the runs reported in Chapter~\ref{ch:results} state which
sampler produced them.

\subsection{Aggregation to Domains}\label{sec:meth-aggregation}

Domain-level rankings pool leaf-level sufficient statistics and refit:
$W_{ij,D_k} = \sum_{L \in \mathcal{D}(D_k)} W_{ij}^{(L)}$ and
analogously for $N_{ij,D_k}$, followed by a fresh MM fit.
Leaf-level parameters are never averaged. Bradley--Terry identification
is fixed only up to a per-component additive constant, so averaging
differently normalised leaf scales produces no interpretable quantity,
whereas pooling counts and refitting yields a single shared scale.

This choice has a consequence that matters for the partition
comparison of Section~\ref{sec:res-c5} and that must be stated where a
reader will see it: under pooling, the partition affects only
\emph{which rubric judged each verdict}, not how the verdicts combine.
The two arms of that comparison are identical in estimation by
construction, so any difference between them is attributable to the
rubric rather than to the roll-up.

\subsection{Validation Against MMLU-Pro}\label{sec:meth-validation}

Ground-truth per-domain accuracy
$A_{i,D_k} = |\{q \in Q_{D_k} : m_i(q) = y_q^{\text{MMLU}}\}| /
|Q_{D_k}|$
is computed strictly after every verdict has been recorded and is used
only for external, post-hoc validation.
Two statistics are reported against it.
The \emph{primary} judge-side statistic is per-verdict agreement
between the judge's winner and the answer key, on the subset where
exactly one response is correct and the judge did not tie. It is a
proportion over thousands of verdicts.
The \emph{secondary} statistic is Spearman $\rho$ between the fitted
$\boldsymbol{\theta}$ and $\mathbf{A}$, which is a rank correlation
over the model count and is therefore heavily quantised. The grid
derivation and the reporting rules that follow from it, including the
requirement that every $\rho$ be reported under both tie conventions,
are in Section~\ref{sec:exp-metrics}.

\subsection{What Is Not Implemented}\label{sec:meth-not-implemented}

Several mechanisms described in earlier design documents for this
system --- per-pair and per-leaf state machines, a global escape
valve, fractional multi-leaf weighting, and five others --- have no
implementation, and none is load-bearing for any result.
Appendix~\ref{app:not-implemented} inventories them in one place, so
that no reader finds a gap the thesis did not name.

\section{Measurement Discipline}\label{sec:meth-discipline}

This project adopted the following rules after violating each of them.
They are stated here, before any result, because Chapter~\ref{ch:results}
reports several nulls and a reader is entitled to know what standard
they were read against. A rule presented as foresight is less credible
than one presented with its scar, so each is given with the instance
that produced it.

\textbf{D1 --- A test that cannot fail is not a test.}
Before running a check, ask what result would falsify it; if the answer
is ``none available'', a clean result carries no information.
Instance: the corpus validator has a \texttt{TRACE\_PRESENCE} check
that measures whether a model's stored responses contain reasoning
text, and that check has \emph{no failure branch at all} --- it is
advisory, and a model with zero stored reasoning is admitted. Six of
the ten models the system hard-excludes by name are excluded for
exactly the condition that check measures and declines to act on; the
resulting policy is a curated list, not a threshold
(Section~\ref{sec:exp-rosters}).\footnote{A second instance: an audit that
classified every identifier column by set membership against every
known identifier space returned approximately $100\%$ everywhere,
because these are dense small-integer ranges over one corpus. The
spaces coincide as sets while disagreeing about which row each value
points to.}

\textbf{D2 --- Report a join's overlap before reporting the join's
result.}
A join with little key overlap produces a clean-looking null
indistinguishable from a real one, so any analysis concluding ``no
effect'' must first report the number of join groups that could have
shown an effect.
Instance, and the one that reached the results: the corpus's
\texttt{mmlu\_pro.id} is not the evaluation table's
\texttt{question\_id}. The numeric join between them succeeds, returns
well-formed rows, and is wrong. It is what collapsed the law domain
from 287 questions to 20 and inverted a domain recommendation
(Section~\ref{sec:res-law}). Joins in this project go through query
text.\footnote{An earlier instance: a design change was made on
measured evidence of ``0 of 36 rescues''. The join was on an iteration
counter that one code path hardcoded to $-1$, so the numerator was
structurally impossible; the mixed-group count was $0$ of $613$ before
the defect was fixed and $97$ of $613$ after --- a single figure that
would have caught it at no cost. The corrected join gives 3 rescues in
404, and a subsequent run refuted the change.}

\textbf{D3 --- Report a median with its interquartile range, and claim
only what the spread supports.}
Instance: $\mathrm{SE}(\Delta J)$ varies by more than an order of
magnitude across sites on this corpus
(Section~\ref{sec:res-construction}), which is the entire reason the
acceptance gate is expressed in standard-error units.
Instance: the acceptance gate's apparent dispersion degradation --- a
coefficient of variation that rose mechanically because the mean fell,
initially read as a real change --- which dissolves under this rule;
the full statistical reading is reported with the result
(Section~\ref{sec:res-construction}).

\textbf{D4 --- A threshold derived for one quantity at one stage does
not transfer.}
Instance: a single separation constant served split acceptance,
component selection, sibling merging and the global proposal gate;
because one of those gates weights its divergence by the fitted
concentration, the same constant meant about $8.1^{\circ}$ of angular
tolerance at freshly split nodes and under $0.7^{\circ}$ at deep
ones, which is why it almost never fired.\footnote{Two further
instances. A second constant served both as a \emph{level} (``is this
partition separated at all?'') and as a \emph{difference} (``does one
more cluster buy enough over the last?''); these are different
quantities and were decoupled into separate fields. A third, a float
tolerance, serves simultaneously as a per-edit acceptance band and as
a per-iteration fixed-point tolerance; the acceptance half has since
been replaced by the standard-error gate, and the certificate half
still borrows a constant derived for the other question.}

\textbf{D5 --- Report a rank correlation under both tie conventions.}
This rule is newer than the other four and was forced by the results of
Section~\ref{sec:res-tie-policy}. Half-weighting ties and dropping tied
comparisons are both defensible, and on the same verdict set they move
Spearman $\rho$ by up to $\pm0.029$ in a direction that is not
consistent across conditions. On one paired contrast the two
conventions disagree about which arm is ahead by an amount larger than
the pre-registered decisive threshold. A single $\rho$ from this
pipeline is therefore not a result; a pair of $\rho$ values under
stated conventions is.

\section{Methodological Assumptions}\label{sec:meth-assumptions}

Table~\ref{tab:meth-assumptions} lists the assumptions the method rests
on, what would go wrong if each failed, and the diagnostic that would
show it.

\begin{table}[htbp]
    \centering
    \caption{Methodological assumptions, failure modes, and diagnostics.}
    \label{tab:meth-assumptions}
    \begin{tabular}{p{3.4cm}p{4.4cm}p{4.4cm}}
        \toprule
        Assumption & Why it might fail & What would detect it \\
        \midrule
        Directional geometry reflects semantic compatibility
        & Embedding direction encodes surface form rather than task
        content
        & The chance-corrected separation score against its two nulls
        (Section~\ref{sec:arch-requirements}) \\
        \addlinespace
        One directional component suffices per node
        & A node holds two distinct sub-topics at comparable mass
        & The min-pair separation gate at split time \\
        \addlinespace
        Bradley--Terry transitivity holds locally within a coherent cell
        & Cyclic dominance survives even inside a leaf
        & Per-cell fit residuals; pairwise-winner accuracy \\
        \addlinespace
        The judge is reliable under domain conditioning
        & Rubric conditioning shifts verdicts without improving them
        & The generic-judge and no-partition arms
        (Section~\ref{sec:res-c5}) \\
        \addlinespace
        Cells carry different true model rankings
        & The corpus is saturated for the roster, so cell identity adds
        nothing to global capability ordering
        & The pre-registered granularity analysis
        (Section~\ref{sec:res-granularity}) \\
        \addlinespace
        Comparisons are independent given the Bradley--Terry parameters
        & Query reuse induces within-query dependence
        & None applied; the Fisher standard errors treat comparisons as
        independent and therefore understate uncertainty \\
        \bottomrule
    \end{tabular}
\end{table}

Three of these are load-bearing, and two of the three did not hold.
The judge-reliability assumption is the one RQ1 is about, and
Section~\ref{sec:res-judge} reports what the judge actually conditions
on. The assumption that cells carry different true rankings is the
premise of the original second research question, and
Section~\ref{sec:res-granularity} reports it as refused under
pre-registration. The independence assumption is known to be violated:
the Fisher standard errors are optimistic by an unmeasured margin, and
the query-level bootstrap intervals of
Section~\ref{sec:exp-metrics} are the reported defence against it.

\section{Persistence and Reproducibility}\label{sec:arch-infra}

State is partitioned across four SQLite stores with distinct write
lifecycles. \texttt{embeddings\_cache.db} is written once before
construction and holds the full-dimensional embedding vectors, the
single source of truth for every downstream slice.
\texttt{mmlu\_pro\_dataset\_cache\_v2.db} holds the MMLU-Pro corpus and
the per-model stored evaluations that supply both the judged responses
and the post-hoc accuracy oracle.
\texttt{snapshots.db} accumulates frozen taxonomy snapshots, each
identified by a UUID; only frozen snapshots are eligible as arena
inputs, and a frozen snapshot is read-only.
\texttt{ratings.db} records every verdict, the per-round estimates, and
run-level metadata, with the snapshot identifier on every row so that
the taxonomy a verdict was produced under is recoverable
unambiguously.

Any result is specified by the tuple
\texttt{(corpusVersion, snapshotId, runId)}, and per-run seeds and
realised call counts are recorded in a manifest emitted with each run.
Two qualifications belong with that claim. The manifest's
working-tree-clean flag does not distinguish a modified source file
from a modified thesis file, so it fires on benign conditions and
should not be read as a reproducibility guarantee on its own. And
\texttt{ratings.db} is keyed by snapshot and condition in a way that
makes a re-run of the same condition \emph{clear} the previous rows
rather than accumulate beside them; the results of
Chapter~\ref{ch:results} were preserved by copying the store before
each subsequent run.
